% !TEX root = ../main.tex
\chapter{Development Process}

The team adopted a \textbf{Scrum-inspired process}, adapted to three developers working part-time on a single codebase.
All members contributed to the design, implementation, and review of the game. In addition, Alex Santini acted as
\textbf{Product Owner}, maintaining and prioritising the product backlog, while Matilde D'Antino acted as
\textbf{Client and Domain Expert}, checking completed increments against the expected game behaviour.

Development was organised into \textbf{four non-consecutive one-week sprints}, each planned for roughly fifteen hours per member,
preceded by an organisational phase that established the build, repository structure, and initial backlog. Each sprint
was expected to end with a \textbf{tangible client-facing increment}, rather than an isolated technical layer.

\section{Task breakdown over time}

The \textbf{product backlog} was drawn up collectively in the first meeting. Work was expressed as \textbf{epics}, each covering one
coherent area of the system, which were then decomposed into smaller \textbf{tasks}.
Only the highest-priority epics were decomposed immediately, while the rest were refined during later
sprint plannings.

Fourteen product epics were identified. The first ten formed the \textbf{core backlog}; the last four were marked
\textbf{optional} from the outset and scheduled only once the core game was playable.

The core epics were:

\begin{enumerate}
    \item Project configuration and CI/CD
    \item Domain model
    \item Map loading, generation and validation
    \item Movement and input management
    \item Game loop and state updates
    \item Collectibles and bonuses
    \item Lives, victory and defeat
    \item Enemy AI standard strategies
    \item User interface and system integration
    \item Documentation
\end{enumerate}

The optional epics were:

\begin{enumerate}[resume]
    \item Scoring and leaderboard
    \item Additional enemy strategy
    \item Game saving and loading
    \item Additional game modes
\end{enumerate}

The \href{https://github.com/PPS-25/PPS-25-scala-man/issues}{backlog is kept on GitHub}, with \textbf{Issues} representing epics
and tasks and the \textbf{GitHub Project board} tracking their state. An item moves through the following workflow:

\begin{center}
  \texttt{backlog} $\rightarrow$ \texttt{ready} $\rightarrow$ \texttt{in progress} $\rightarrow$
  \texttt{in review} $\rightarrow$ \texttt{done}
\end{center}

Here, \texttt{backlog} identifies work not yet scheduled into a sprint, \texttt{ready} scheduled but unstarted work,
and \texttt{in review} work with an open pull request.

\section{Planned meetings and interactions}
Two recurring meetings were planned for each sprint:

\begin{itemize}
  \item \textbf{Friday backlog refinement}: the team re-examined current work, splitting unexpectedly complex items
  into smaller tasks that could move to the following sprint.
  \item \textbf{Sunday sprint boundary}: the team held the sprint review, retrospective, and planning for the following
  week; the completed increment was demonstrated and judged against the definition of done.
\end{itemize}

Outside these fixed points, the team interacted continuously through pull request reviews.

The planned sprint windows were:
\begin{center}
    \begin{tabular}{|l|l|}
        \hline
        \textbf{Sprint} & \textbf{Period} \\
        \hline
        1 & 3-9 August \\
        2 & 10-16 August \\
        3 & 24-30 August \\
        4 & 7-13 September \\
        \hline
    \end{tabular}
\end{center}

The following \textbf{increments} were concluded across the four sprints:

\begin{enumerate}
  \item \textbf{The project foundation}: build, automated checks, the core domain, and validated textual maps.
  \item \textbf{The gameplay core}: movement, collisions, lives, effects, enemy strategies, and scoring foundations.
  \item \textbf{A playable ScalaFX application}: rendering, menu, commands, saves, leaderboards, and game modes.
  \item \textbf{The final delivery}: Patroller integration, refinements, documentation, regression checks, and packaging.
\end{enumerate}

\section{Task review over time}

Work was reviewed continuously, through the \textbf{branching model} and the \textbf{pull request process}, rather than only at sprint
boundaries. The main development flow used four kinds of branch:

\begin{center}
    \begin{tabular}{|l|p{9cm}|}
        \hline
        \textbf{Branch} & \textbf{Purpose} \\
        \hline
        \texttt{main} & Release branch, kept in a state that can always be built and released \\
        \texttt{dev} & Integration branch, where completed epics are brought together \\
        \texttt{epic/<n>-<name>} & Branch collecting all work belonging to a single epic \\
        \texttt{task/<n>-<name>} & Branch for an individual task within an epic \\
        \hline
    \end{tabular}
\end{center}

Work flows in one direction:

\begin{center}
\texttt{task} $\rightarrow$ \texttt{epic} $\rightarrow$ \texttt{dev} $\rightarrow$ \texttt{main}
\end{center}

Short-lived \texttt{docs/}, \texttt{refactor/}, \texttt{chore/}, and \texttt{integrate/} branches complemented this
flow for documentation, maintenance, and conflict resolution during the final review phase.

The team used \textbf{pull requests} to integrate work into \texttt{dev}, with reviews performed by another team member before
merging. This provided an opportunity to discuss changes and identify potential integration issues before they reached
the shared development branch.

Commit messages follow the \textbf{Conventional Commits} format, \texttt{<type>(<scope>): description}, which keeps the history
readable.

\textbf{Definition of done.} A task is done when it:

\begin{itemize}
  \item fulfils its \textbf{requirements};
  \item is covered by \textbf{automated tests};
  \item passes \textbf{continuous integration}; and
  \item has its \textbf{documentation} updated.
\end{itemize}

An epic is done when all of its tasks are done and its work has been merged into \texttt{dev}.

\section{Testing, build, and CI tools}

The project is written in \textbf{Scala 3.3.5} and is built and tested with \textbf{JDK 21}. The build is managed by \textbf{sbt}, which handles
compilation, dependency resolution, test execution and packaging through a single build definition. Two library dependencies
are declared: \textbf{ScalaFX 21.0.0-R32} for the graphical interface, and \textbf{ScalaTest 3.2.18}, scoped to testing only.

\subsection{Code style}

Formatting is delegated to \textbf{scalafmt}, driven by the sbt-scalafmt plugin and configured in \texttt{.scalafmt.conf}.
Continuous integration runs \texttt{scalafmtCheckAll} and fails on formatting deviations, enforcing the agreed style
mechanically and reducing formatting-related discussion during code review.

\subsection{Static analysis}

\textbf{WartRemover} complements formatting by checking production sources for selected unsafe constructs: partial
\texttt{Option} and collection operations, \texttt{throw}, and \texttt{Any}. The checks are deliberately focused on
production code, where they prevent avoidable failures from reaching the game, while test fixtures remain free to use
the concise setup that makes each test readable.

\subsection{Testing}

\textbf{Automated tests} are written with \textbf{ScalaTest} and run under sbt. Contributors are asked to run the suite locally before
pushing, following the documented pre-push checklist; continuous integration runs it again on every push and pull request.

\subsection{Continuous integration}

\textbf{Continuous integration} runs on \textbf{GitHub Actions} on every push and pull request. It uses Temurin JDK 21 and an sbt cache
to compile the project, run the ScalaTest suite, and enforce formatting with scalafmt; a failing check is reported before
the change is integrated.

\subsection{Continuous Delivery}

The \textbf{Project Delivery} workflow runs on every push to \texttt{main} and can also be triggered manually. It builds
an executable \textbf{fat JAR} and compiles the \textbf{report PDF} in a TeX Live container. Both artifacts and the sources are packaged
into a single archive, whose structure is checked before the archive is published as a \textbf{GitHub release}.
